\begin{proofbox}
	\textbf{\textcolor{CProof}{思路提示}}

	利用平行线性质得到对应角相等，利用平行线分线段成比例定理得到对应边成比例，然后直接根据相似三角形定义得证。

	\medskip
	\textbf{\textcolor{CProof}{正式推导}}
	\begin{description}[style=nextline, leftmargin=2.8em, labelsep=0.5em, itemsep=0.55em]
		\item[\textbf{步骤一（找等角）：}] 因为 $DE\parallel BC$，由平行线性质，同位角相等，得
			\[
				\angle ADE = \angle ABC,\quad \angle AED = \angle ACB.
			\]
			又 $\angle A = \angle A$（公共角），所以两三角形对应角全部相等。
			\par\textbf{理由：} 平行线同位角相等，公共角相同。

		\item[\textbf{步骤二（比例关系）：}] 由平行线分线段成比例定理（平行于三角形一边的直线截其他两边，所得线段对应成比例），有
			\[
				\begin{aligned}
					\frac{AD}{AB} &= \frac{AE}{AC} \\
					&= \frac{DE}{BC}.
			\end{aligned}
			\]
			\par\textbf{理由：} 定理直接应用。

		\item[\textbf{步骤三（定义证相似）：}] 根据相似三角形定义：对应角相等且对应边成比例的两个三角形相似，故
			\[
				\triangle ADE \sim \triangle ABC.
			\]
			\par\textbf{理由：} 步骤一已证角相等，步骤二已证边成比例。
	\end{description}

	\medskip
	\textbf{\textcolor{CProof}{结论回扣}}

	\textbf{上述推导证实了“平行于三角形一边的直线截得的三角形与原三角形相似”这一判定。只要$DE\parallel BC$，就有$\triangle ADE\sim\triangle ABC$。}
\end{proofbox}
